\begin{textsample}{NEG Dim 4 – 2015\_c5 – Score 1.00 – 2015\_c5/enade\_2015\_tec\_gest\_rec\_hum\_77f5a8062b.txt}  \label{ex:f4_neg_183}
CAPÍTULO 5 DISTRIBUIÇÃO DOS CONCEITOS 5.1 PANORAMA NACIONAL DA DISTRIBUIÇÃO DOS CONCEITOS O Conceito Enade é um indicador de qualidade que avalia o desempenho dos estudantes a partir dos resultados do Enade.
Os valores possíveis vão de 1 ( pior situação ) a 5 ( melhor situação ).
As Unidades de Observação com apenas um ou sem nenhum concluinte participante não obtêm o Conceito Enade, ficando Sem Conceito ( SC ).
A Tabela 5.1 apresenta a quantidade e distribuição de cursos de Tecnologia em Gestão de Recursos Humanos participantes do Enade/2015, por faixa de conceito e Grande Região.
Mais uma vez, ressaltando que a diferença entre os cursos tabulados neste capítulo e no capítulo 2 são os cursos sem conceito, em princípio, aqueles sem alunos concluintes que participassem da prova ou que tivessem tão somente um aluno concluinte18.
Observando -se os dados da Tabela 5.1, nota -se que, dos 523 cursos participantes, 242 ( 46,3 \% ) classificaram -se com conceito 3, o valor modal.
Este foi também o conceito modal em quase todas regiões, exceto na região Sul.
O conceito 4 foi o segundo mais frequente em nível nacional ( 24,5 \%, correspondendo a 128 cursos ) e o conceito 2, o terceiro ( 20,8 \%, correspondendo a 109 cursos ).
Houve, ainda, 22 cursos ( 4,2 \% ) que receberam conceito 5 e seis ( 1,1 \% ) cursos que receberam conceito 1.
Dos 523 cursos de Tecnologia em Gestão de Recursos Humanos, 16 cursos ( 3,1 \% ) ficaram sem conceito ( SC ). 18 Estes cursos com somente um concluinte não foram considerados no capítulo 2.
A região Norte participou com 17 cursos ou 3,3 \% do total nacional.
Nesta região, o conceito 3 foi atribuído a 11 cursos dos 17 participantes, o que equivale a 64,7 \% do total regional, além de ser o conceito modal na região.
Um curso ficou sem conceito ( SC ) e nenhum recebeu os conceitos 1 ou 5.
Três cursos ( 17,6 \% ) foram avaliados com conceito 2 e dois cursos ( 11,8 \% ) com conceito 4.
A região Nordeste participou com 69 cursos ou 13,2 \% do total nacional.
Nesta região, 34 cursos ( 49,3 \% em termos regionais ) obtiveram conceito 3, o conceito modal para a região.
Os conceitos 2 e 4 foram atribuídos a, respectivamente, 17 e 11 cursos ( 24,6 \% e 15,9 \% ).
Um curso ( 1,4 \% ) obteve o conceito 1 e dois ( 2,9 \% ) o conceito 5.
Quatro cursos ficaram sem conceito ( SC ) nessa região.
Dos 314 cursos participantes da região Sudeste ( 60,0 \% do total nacional ), o conceito 3 também foi o conceito modal, e foi obtido por 148 cursos ( 47,1 \% ).
O conceito 1 foi atribuído a dois cursos ( 0,6 \% ) e os conceitos 2 e 4, a 68 e 75 cursos, correspondendo a 21,7 \% e 23,9 \%, respectivamente.
O conceito 5 foi atribuído a 12 cursos ( 3,8 \% ).
Nove cursos ( 2,9 \% ) ficaram sem conceito ( SC ).
Os 82 cursos da região Sul ficaram distribuídos em todas as faixas de conceitos.
Há predominância do conceito 4 com 40,2 \% ( conceito modal ), ou seja, 33 dos 82 cursos participantes na região Sul.
O conceito 2 foi atribuído a 13 cursos ( 15,9 \% ) e o conceito 3, a 27 cursos ( 32,9 \% ).
O conceito 5 foi recebido por seis cursos ( 7,3 \% ) e o conceito 1 por um curso ( 1,2 \% ).
Dois cursos na região Sul ficaram sem conceito ( SC ).
Já dos 41 cursos participantes na região Centro-Oeste, 22 cursos ( 53,7 \% ) receberam conceito 3, o conceito modal.
Nesta região, nenhum curso ficou sem conceito ( SC ).
Já oito cursos foram \textbf{alocados} no conceito 2, correspondendo a 19,5 \% dos cursos nessa região.
Sete cursos foram avaliados com conceito 4 ( 17,1 \% ), e quatro cursos ficaram com conceito 1 e 5 ( dois cursos em cada conceito ), correspondendo 4,9 \%, cada.
O Gráfico 5.1 apresenta a distribuição do Conceito Enade segundo Grande Região : poligonais mais à esquerda representam Grandes Regiões com pior distribuição, e poligonais mais à direita, Grandes Regiões com melhores conceitos.
A região Sul ( linha roxa ) apresenta os cursos em IES com conceitos mais altos, ao passo que as outras região, Norte ( linha azul ), Sudeste ( linha vermelha ), Nordeste ( linha verde ) e Centro-Oeste ( linha preta ) apresentam linhas à esquerda e que se entrecruzam. 5.2 CONCEITOS POR CATEGORIA ADMINISTRATIVA E POR GRANDE REGIÃO A Tabela 5.2 apresenta a distribuição dos cursos participantes do Enade/2015 de Tecnologia em Gestão de Recursos Humanos, por Categoria Administrativa, de acordo com os conceitos por eles alcançados, segundo as Grandes Regiões brasileiras.
Dos 523 cursos participantes, oito ( 1,5 \% ) eram ministrados em Instituições Públicas e 515 ( 98,5 \% ), em Instituições Privadas.
De acordo com as informações da Tabela 5.2, em termos nacionais, entre Instituições Públicas e Privadas, das 22 IES que receberam o conceito 5, duas eram Públicas e 20 Privadas.
Dos oito cursos participantes de IES Públicas, o conceito 4 foi o valor modal, atribuído a três cursos ( 37,5 \% ).
Entre os demais cursos de Instituições Públicas participantes, quatro cursos foram avaliados com conceitos 3 e 5 ( dois cursos em cada conceito ).
Nesta categoria, somente um curso ( 12,5 \% ) ficou sem conceito ( SC ).
Além disso, nenhum curso obteve os conceitos 1 ou 2.
Na rede privada, o conceito modal foi 3, com 240 cursos dos 515 da categoria.
Entre os demais cursos participantes, seis receberam conceito 1, e 109, conceito 2.
O conceito 4 foi atribuído a 125 cursos e o conceito 5, a 20 cursos.
Nesta Categoria Administrativa, outros 15 cursos ficaram sem conceito ( SC ).
Na análise por região, observa -se que, na região Norte, as Instituições Privadas participaram com todos os 17 cursos nesta região, dos quais um ficou sem conceito ( SC ) e nenhum com conceitos 1 ou 5.
O conceito modal para as Instituições Privadas na região foi 3, com 11 cursos, correspondendo a 64,7 \%.
Três cursos foram \textbf{alocados} no conceito 2 e dois, no conceito 4.
As Instituições Públicas não tiveram participação nesta região.
Na região Nordeste, a rede privada concentrou 68 dos 69 cursos participantes, o equivalente a 98,6 \% do total da região.
Dentre os cursos oferecidos por IES Privadas no Nordeste, um curso ficou \textbf{alocado} ao conceito 1, o conceito 2 foi atribuído a 17 cursos, 34 cursos receberam o conceito 3 ( conceito modal ), dez, o conceito 4 e dois, o conceito 5.
Além disso, quatro cursos não obtiveram conceito ( SC ).
As Instituições Públicas dessa região participaram com somente um curso, e este foi \textbf{alocado} ao conceito 4.
Na região Sudeste, a proporção de cursos da rede privada, 97,8 \%, foi menos elevada do que nas demais regiões brasileiras, correspondendo a 307 dos 314 cursos participantes.
Nesta categoria, na região Sudeste, o conceito modal foi 3 ( 146 cursos ).
Nesta combinação de Categoria Administrativa e Grande Região, oito cursos não receberam conceito ( SC ).
Os demais foram avaliados com conceito 1 ( dois cursos ), conceito 2 ( 68 cursos ), conceito 4 ( 73 cursos ) e conceito 5 ( dez cursos ).
Entre os sete cursos oferecidos por Instituições Públicas na região Sudeste, seis foram \textbf{alocados} nos conceitos 3, 4 e 5 ( dois em cada conceito ).
Na região Sudeste, nesta categoria, somente um curso ficou sem conceito ( SC ) e nenhum recebeu os conceitos 1 ou 2.
As Instituições Privadas concentraram todos os cursos participantes da região Sul, correspondendo a 82 cursos.
Desses, 33 ficaram com conceito 4, o conceito modal.
Nesta combinação de Categoria Administrativa e Grande Região, 27 cursos ficaram com conceito 3, 13 cursos obtiveram conceito 2, enquanto Seis cursos receberam conceito 5.
Um único curso foi associado ao conceito 1, enquanto dois cursos ficaram sem conceito ( SC ).
As Instituições Públicas na região Sul não tiveram participação.
Na região Centro-Oeste, todos os cursos participantes eram de Instituições Privadas equivalente a 41 cursos.
Destes, 22 concentraram -se no conceito 3, conceito modal.
Dentre os demais, oito receberam conceito 2, e sete receberam conceito 4.
Nenhum curso ficou sem conceito ( SC ) e quatro receberam os conceitos 1 e 5 ( dois cursos em cada conceito ).
As Instituições Públicas também não tiveram participação nesta região.
O Gráfico 5.2 apresenta a distribuição do Conceito Enade segundo a Categoria Administrativa da IES.
Os cursos em IES Públicas ( linha azul ) apresentam uma poligonal mais à direita do que os em IES Privadas ( linha verde ) e, consequentemente, uma distribuição de Conceitos Enade com valores maiores.
As Instituições Privadas participaram com todos os 17 cursos nesta região, dos quais um ficou sem conceito ( SC ) e nenhum com conceitos 1 ou 5.
O conceito modal para as Instituições Privadas na região foi 3, com 11 cursos, correspondendo a 64,7 \%.
Três cursos foram \textbf{alocados} no conceito 2 e dois, no conceito 4.
As Instituições Públicas não tiveram participação nesta região.
Na região Nordeste, a rede privada concentrou 68 dos 69 cursos participantes, o equivalente a 98,6 \% do total da região.
Dentre os cursos oferecidos por IES Privadas no Nordeste, um curso ficou \textbf{alocado} ao conceito 1, o conceito 2 foi atribuído a 17 cursos, 34 cursos receberam o conceito 3 ( conceito modal ), dez, o conceito 4 e dois, o conceito 5.
Além disso, quatro cursos não obtiveram conceito ( SC ).
As Instituições Públicas dessa região participaram com somente um curso, e este foi \textbf{alocado} ao conceito 4.
Na região Sudeste, a proporção de cursos da rede privada, 97,8 \%, foi menos elevada do que nas demais regiões brasileiras, correspondendo a 307 dos 314 cursos participantes.
Nesta categoria, na região Sudeste, o conceito modal foi 3 ( 146 cursos ).
Nesta combinação de Categoria Administrativa e Grande Região, oito cursos não receberam conceito ( SC ).
Os demais foram avaliados com conceito 1 ( dois cursos ), conceito 2 ( 68 cursos ), conceito 4 ( 73 cursos ) e conceito 5 ( dez cursos ).
Entre os sete cursos oferecidos por Instituições Públicas na região Sudeste, seis foram \textbf{alocados} nos conceitos 3, 4 e 5 ( dois em cada conceito ).
Na região Sudeste, nesta categoria, somente um curso ficou sem conceito ( SC ) e nenhum recebeu os conceitos 1 ou 2.
As Instituições Privadas concentraram todos os cursos participantes da região Sul, correspondendo a 82 cursos.
Desses, 33 ficaram com conceito 4, o conceito modal.
Nesta combinação de Categoria Administrativa e Grande Região, 27 cursos ficaram com conceito 3, 13 cursos obtiveram conceito 2, enquanto Seis cursos receberam conceito 5.
Um único curso foi associado ao conceito 1, enquanto dois cursos ficaram sem conceito ( SC ).
As Instituições Públicas na região Sul não tiveram participação.
Na região Centro-Oeste, todos os cursos participantes eram de Instituições Privadas equivalente a 41 cursos.
Destes, 22 concentraram -se no conceito 3, conceito modal.
Dentre os demais, oito receberam conceito 2, e sete receberam conceito 4.
Nenhum curso ficou sem conceito ( SC ) e quatro receberam os conceitos 1 e 5 ( dois cursos em cada conceito ).
As Instituições Públicas também não tiveram participação nesta região.
O Gráfico 5.2 apresenta a distribuição do Conceito Enade segundo a Categoria Administrativa da IES.
Os cursos em IES Públicas ( linha azul ) apresentam uma poligonal mais à direita do que os em IES Privadas ( linha verde ) e, consequentemente, uma distribuição de Conceitos Enade com valores maiores. 5.3 CONCEITOS POR ORGANIZAÇÃO ACADÊMICA E POR GRANDE REGIÃO Na Tabela 5.3, encontra -se a distribuição dos conceitos atribuídos aos cursos participantes do Enade/2015 na Área de Tecnologia em Gestão de Recursos Humanos, por Organização Acadêmica, segundo as Grandes Regiões brasileiras.
Dos 523 cursos de Tecnologia em Gestão de Recursos Humanos participantes, 151 eram oferecidos em Universidades, 106 em Centros Universitários e 266 em Faculdades.
Esta distribuição corresponde a, respectivamente, 28,9 \%, 20,3 \% e 50,9 \% dos cursos.
De acordo com os dados apresentados, de todos os 22 cursos avaliados com conceito 5, dez eram vinculados a Universidades.
Seis cursos oferecidos em Universidades ficaram sem conceito ( SC ).
Esse tipo de Organização Acadêmica teve o conceito 3 como modal, com 69 cursos ( 45,7 \% ).
Os demais cursos de Universidades avaliados receberam os conceitos 1 ( dois cursos ), 2 ( 27 cursos ), 4 ( 37 cursos ) e conceito 5 ( dez cursos, como já mencionado ).
Entre os cursos em Centros Universitários, o conceito modal foi 3, com 50 cursos ( 47,2 \% ).
Neste tipo de Organização Acadêmica, nenhum curso ficou sem conceito ( SC ), e os outros cursos receberam os conceitos 1 ( um curso ), 2 ( 21 cursos ), 4 ( 30 cursos ) e 5 ( quatro cursos ).
Nas Faculdades, dez dos 266 cursos ficaram sem conceito ( SC ), e 123 ( 46,2 \% ) ficaram \textbf{alocados} no conceito 3 ( conceito modal ), seguido de 122 cursos que obtiveram conceitos 2 e 4 ( 61 cursos em cada conceito ).
Os demais cursos neste tipo de Organização Acadêmica, receberam os conceitos 1 ( três cursos ) e 5 ( oito cursos ).
Considerando -se separadamente as regiões brasileiras, verifica -se que, na região Norte, as Universidades concentraram apenas um dos 17 cursos participantes, e este foi \textbf{alocado} no conceito 4.
Os Centros Universitários da região Norte foram representados por dois cursos, e estes receberam os conceitos 2 e 3 ( um curso em cada conceito ).
As Faculdades participaram com 14 cursos na região Norte que receberam conceitos 3 ( dez cursos, conceito modal ), conceito 2 ( dois cursos ) e 4 ( um curso ).
Além disso, um curso ficou sem conceito ( SC ).
Na região Nordeste, as Universidades participaram com dez dos 69 cursos na Área de Tecnologia em Gestão de Recursos Humanos.
O conceito modal foi 3, com sete cursos, enquanto o conceito 3 foi atribuído a dois cursos e o conceito 4 a um curso.
Os Centros Universitários contaram com 11 cursos participante na região Nordeste.
Destes, quatro obtiveram conceito 4, conceito modal.
Além disso, seis cursos foram \textbf{alocados} nos conceitos 2 e 3 ( três cursos em cada conceito ) e um curso obteve conceito 5.
As Faculdades foram representadas por 48 cursos na região Nordeste, dos quais 24 receberam conceito 3 ( conceito modal ), 12 ficaram com conceito 2 e seis com conceito 4.
Dois cursos ficaram \textbf{alocados} nos conceitos 1 e 5 ( um curso em cada conceito ) e quatro cursos ficaram sem conceito ( SC ).
Na região Sudeste, as Universidades concentraram 114 dos 314 cursos da região.
Entre os cursos em Universidades, nessa região, o conceito modal foi 3 com 53 cursos, enquanto seis cursos ficaram sem conceito ( SC ).
Os demais cursos receberam os conceitos 1 ( um curso ), 2 ( 20 cursos ), 4 ( 27 cursos ) e 5 ( sete cursos ).
Os Centros Universitários participaram com 67 cursos na região Sudeste, dos quais 36 obtiveram o conceito modal, 3, enquanto nenhum curso ficou sem conceito ( SC ) ou recebeu conceito 1.
Os demais receberam os conceitos 2 ( 13 cursos ), 4 ( 17 cursos ) e 5 ( um curso ).
As Faculdades foram representadas por 133 cursos na região Sudeste dos quais 59 receberam o conceito 3 ( conceito modal ).
Os demais se distribuíram nos conceitos 1 ( um curso ), 2 ( 35 cursos ), 4 ( 31 cursos ) e 5 ( quatro cursos ).
Três cursos ficaram sem conceito ( SC ) nessa categoria e região.
Dos 82 cursos da região Sul, 20 eram de Universidades, para os quais os conceitos modais foram 3 e 4, com sete cursos, cada.
Nesse tipo de organização, nenhum dos cursos ficou sem conceito ( SC ) ou com conceito 1.
Os demais receberam os conceitos 2 ( quatro cursos ) e 5 ( dois cursos ).
Dos 15 cursos participantes de Centros Universitários da região Sul, seis obtiveram conceito 4, conceito modal.
Os demais se distribuíram nos conceitos 1 ( um curso ), 2 ( dois cursos ), 3 ( cinco cursos ) e 5 ( um curso ).
Nenhum curso ficou sem conceito ( SC ).
Foram 47 os cursos vinculados a Faculdades na região Sul, e 20 desses receberam conceito 4, o conceito modal.
Sete cursos receberam o conceito 2, enquanto 15 receberam conceito 3 e três, conceito 5.
Dois cursos ficaram sem conceito ( SC ) e nenhum recebeu o conceito 1.
Na região Centro-Oeste, seis dos 41 cursos eram de Universidades.
Nesse tipo de organização, nenhum curso ficou sem conceito ( SC ).
O conceito modal foi 3, com dois cursos.
Os outros quatro cursos ficaram \textbf{alocados} nos demais conceitos : 1, 2, 4 e 5 ( um curso em cada ).
Os Centros Universitários da região Centro-Oeste contaram com 11 cursos, dos quais cinco obtiveram conceito 3 ( conceito modal ), três receberam conceito 4 e dois, conceito 2.
Além disso, nenhum curso ficou sem conceito ( SC ) ou com conceito 1, e um único curso foi \textbf{alocado} ao conceito 5.
Dos 24 cursos oferecidos em Faculdades na região Centro-Oeste, 15 ficaram \textbf{alocados} no conceito 3, conceito modal.
Cinco cursos ficaram com conceito 2, e três, com conceito 4.
Além disso, nenhum curso ficou sem conceito ( SC ) ou com conceito 5 e um único curso foi \textbf{alocado} no conceito 1.
O Gráfico 5.3 apresenta a distribuição do Conceito Enade segundo a Organização Acadêmica da IES.
Os cursos em Universidades ( linha azul ) e Centros Universitários ( linha verde ) apresentam uma poligonal mais à direita do que as Faculdades ( linha vermelha ), denotando uma distribuição de Conceitos Enade maiores.

% matched lemmas: alocar
\end{textsample}
